\chapter*{空间直线与平面}
\section*{平面及其基本性质}
\subsection*{平面及其基本性质 (1)}

A 组：
\begin{enumerate}
    \item 几何中平面的特征:平面没有\blkda{薄厚}，是\blkda{无限延展}的.
    
    \item 判断下列命题是否正确，正确的打 “ \ding{51} ”，错误的打 “ \ding{55} ”
    \begin{enumerate}
        \item 平行四边形就是平面. \blkda{\ding{55}}
        
        \item 一条直线把它所在的平面分成两部分，一个平面把空间分成两部分. \blkda{\ding{51}}
        
        \item 若平面 $a$ 的面积是 ${20}{m}^{2}$ ，平面 $\beta$ 的面积是 ${10}{m}^{2}$ ，则平面 $a$ 的面积是平面 $\beta$ 的面积的 2 倍 \blkda{\ding{55}}
        
        \item 五个平面重叠在一起比两个平面重叠在一起要厚. \blkda{\ding{55}}
        
        \item 如果一条直线上所有的点都在某一个面上, 那么这个面一定是平面. \blkda{\ding{55}}
        
        \item 经过面内任意两点的直线，若直线上各点都在这个面内，那么这个面是平面. \blkda{\ding{51}}
    \end{enumerate}

    \item 如图所示，用集合的语言表述下列语句:
    
    \begin{center}
    \includegraphics[width=0.2\textwidth]{bodyxb/src/bo_d4g0j5k601uc73fkjrb0_0_1078_1165_323_175_0.jpg}
    \end{center} 
    \begin{dingautolist}{172}
        \item 点 $A\text{ 、 }B$ 在直线 $a$ 上:\blkda[4]{$A\in a , B \in a$}
        
        \item 直线 $a$ 在平面 $\alpha$ 内:\blkda[4]{$a \subset \alpha$}
        
        \item 点 $C$ 在平面 $\alpha$ 内:\blkda[4]{$C \in \alpha$}

        \item 点 $D$ 不在平面 $\alpha$ 内:\blkda[4]{$D \notin \alpha$}
        
        \item 直线 $b$ 不在平面 $\alpha$ 内:\blkda[4]{$b \not \subset \alpha$}.
    \end{dingautolist}
\end{enumerate}

B 组：
\begin{enumerate}
    \item 四个平面将空间最少分成 $m$ 部分，最多分成 $n$ 部分，则 $m + n =$ \blkda{20}.
    \begin{jieda}
        $m = 15, n = 5$
    \end{jieda}
    
    \item 不都在一条直线上的四点可以确定\blkda{1或4}个平面.
    \begin{jieda}
        不存在三个点共线时有4个面（正四面体）

        存在三个点共线时根据公理2的推论1可知有1个面
    \end{jieda}

    \item 一条直线和直线外三点可以确定\blkda[2]{1或3或4}个平面.
    \begin{jieda}
        三点共线时可以确定1个面或者3个面

        三个点不共线时根据公理2可知这三个点确定一个平面$\alpha$，若直线在$\alpha$内则只确定一个平面，否则，直线和三个点会确定$\beta_1, \beta_2, \beta_3$三个平面，因此最多确定4个平面，若直线和某两点的连线有交点，则会使得$\beta_i$与$\beta_j$重合，退化为3个平面.
    \end{jieda}

    \item 两两相交于一点的三条直线可以确定\blkda{1或3}个平面.
    \begin{jieda}
        两条相交直线确定一个平面，若第三个直线不在该面内，则确定三个平面，否则确定一个平面
    \end{jieda}

    \item 两两平行的四条直线可以确定\blkda[2]{1或4或6}个平面.
    \begin{jieda}
        四条直线共面则可以确定一个平面

        三条直线共面，另一条不在面内则确定四个平面

        任意三条直线都不共面则确定六个平面
    \end{jieda}

    \item 由正方体的一条棱和一个不在此棱上的顶点确定的平面有\blkda{12}个.

    \item 下列判断中: 
    \begin{dingautolist}{172}
        \item 三点确定一个平面
        \item 一条直线和一点确定一个平面
        \item 两条直线确定一个平面
        \item 三角形和梯形一定是平面图形
        \item 四边形一定是平面图形
        \item 六边形一定是平面图形
        \item 两两相交的三条直线确定一个平面
    \end{dingautolist}
    其中正确的是\blkda{\ding{175}}.
    
    \item 空间三个平面可以把空间分成\dotfill{(\kh{B})}
    \xx{4 或 6 或 7 个部分}{4 或 6 或 7 或 8 个部分}{4 或 7 或 8 个部分}{6或 7 或 8 个部分}
    
    \item 若直线上有两个点在平面外, 则\dotfill{(\kh{D})}
    \xx{直线上至少有一个点在平面内}{直线上有无穷多个点在平面内}{直线上所有点都在平面外}{直线上至多有一个点在平面内}
    
    \item 已知命题 “直线 $l$ 上两点 $A\text{ 、 }B$ 在平面 $a$ 上”,则下列命题中与其不等价的是\dotfill{(\kh{C})}
    \xx{$l \subset  a$}{平面 $a$ 过直线 $l$}{直线 $l$ 上只有两点在平面 $a$ 上}{直线 $l$ 上所有点都在平面 $a$ 上}
    
    \item 若点 $M$ 在直线 $b$ 上， $b$ 在平面 $\beta$ 上，则 $M$ 、 $b$ 、 $\beta$ 之间的关系可记作\dotfill{(\kh{C})}
    \xx{$M \in  b \in  \beta$}{$M \subset  b \subset  \beta$}{$M \in  b \subset  \beta$}{$M \subset  b \in  \beta$}
    
    \item 请将以下四图中，看得见的部分用实线描出.
    \fourtu{\includegraphics[width=0.6\textwidth]{bodyxb/src/bo_d4g0j5k601uc73fkjrb0_1_258_1352_241_197_0.jpg}}{\includegraphics[width=0.6\textwidth]{bodyxb/src/bo_d4g0j5k601uc73fkjrb0_1_557_1363_304_177_0.jpg}}{\includegraphics[width=0.6\textwidth]{bodyxb/src/bo_d4g0j5k601uc73fkjrb0_1_916_1355_247_199_0.jpg}}{\includegraphics[width=0.6\textwidth]{bodyxb/src/bo_d4g0j5k601uc73fkjrb0_1_1227_1357_225_194_0.jpg}}
    
    \item 画图表示下列集合语言所描述的关系:
    \begin{enumerate}
        \item $A \in  \alpha ,a \subset  \alpha ,A \notin  a$
        \item $a \cap  \beta  = 1,A \in  \alpha$ ，且 $A \in  \beta$
        \item $a,b,c \subset  \alpha$ ,且 $a \cap  b = A,b \cap  c = B,c \cap  a = C$
        \item $A \in  \alpha, A \in  l,l \cap  \beta  = B, \alpha \cap  \beta  = m,B \in  m$
        \item $a \cap  \beta  = \varnothing ,l \cap  a = A,l \cap  \beta  = B$ 
        \item ${l}_{1} \subset  \alpha,{l}_{2} \cap  \alpha = M,M \notin  {l}_{1}$
    \end{enumerate}
    
    \item 已知直线 $a$ 和直线 $b$ 相交于点 $M$ ，点 $M$ 在平面 $\alpha$ 上，试用集合语言描述上述位置关系, 并画出图形.
    \begin{jieda}
        $a \cap b = M, M \in \alpha$
    \end{jieda}
    
    \item 已知 $I$ 为 $\bigtriangleup  {ABC}$ 的内心. 若 $A$ 、 $B$ 、 $C \in  \alpha$ ，利用公理 1 证明: $I \in  \alpha$ .
    \begin{jieda}
        连接$AI$交$BC$于$M$，根据公理1可知$M \in \alpha$，故$AM \subset \alpha$，因此$I \in \alpha$
    \end{jieda}
    
    \item 用反证法证明: 一个平面 $\alpha$ 与不在 $\alpha$ 上的一条直线 $l$ 最多只能有一个交点.
    \begin{jieda}
        假设$\alpha$与$l$存在两个及以上的公共点，设其中两个为$A,B$，则$A, B \in \alpha$，根据公理1可知$l \subset \alpha$，矛盾，因此一个平面 $\alpha$ 与不在 $\alpha$ 上的一条直线 $l$ 最多只能有一个交点.
    \end{jieda}
\end{enumerate}

\clearpage
\subsection*{平面及其基本性质 (2)}

A 组：
\begin{enumerate}
    \item 判断下列命题是否正确，正确的打 “ \ding{51} ”，错误的打 “ \ding{55} ”
    \begin{dingautolist}{172}
        \item 若 $\bigtriangleup  {ABC}$ 的两边 ${AB}$ 、 ${AC}$ 在平面 $\alpha$ 上，则它的第三边 ${BC}$ 也在平面 $a$ 上\blkda{\ding{51}}
        \item 若两个平面有不同的三个公共点，则这两个平面重合\blkda{\ding{55}}
        \item 若四点不共面，则每三个点一定不共线\blkda{\ding{51}}
        \item 已知三个平面两两相交，且有三条交线. 如果这三条交线中有两条交线平行，那么第三条交线也与另两条直线平行\blkda{\ding{51}}
        \item 四边相等的四边形是平面图形\blkda{\ding{55}}
    \end{dingautolist}
    
    \item 已知 $a \cap  \beta  = l,m \subset  a,n \subset  \beta ,m \cap  n = P$ ，则点 $P$ 与直线 $l$ 的位置关系用相应的符号表示为\blkda{$P \in l$}.
    \begin{jieda}
        两平面的公共点都在交线上.
    \end{jieda}
\end{enumerate}

B 组：
\begin{enumerate}
    \item 下列命题
    \begin{dingautolist}{172}
        \item 若空间四点共面，则其中必有三点共线
        \item 若空间四点中有三点共线，则此四点必共面
        \item 若空间四点中任何三点不共线，则此四点不共面
        \item 若空间四点不共面，则其中任意三点不共线
    \end{dingautolist}
    其中真命题的个数是\dotfill{(\kh{B})}
    \xx{1 个}{2 个}{3 个}{4 个}
    \begin{jieda}
        \begin{dingautolist}{172}
            \item 四点共面不一定得到三点共线，比如平面四边形，故错误
            \item 若空间四点中有三点共线，则此四点必共面，故正确
            \item 若空间四点中任何三点不共线，则此四点可能共面，比如平面四边形，故错误
            \item 若空间四点不共面，则其中任意三点不共线；假设其中三个点共线，则第四个点要么与这三点在一条直线上，要么在直线外，根据直线和直线外一点可确定一个平面可知，这四点共面，矛盾，故任意三点不共线；故正确
        \end{dingautolist}
    \end{jieda}

    \item 在正方体中， ${E}$ 、 ${F}$ 、 ${G}$ 、 ${H}$ 分别是该点所在棱的中点，则下列图形中 ${E}$ 、 ${F}$ 、 ${G}$ 、 ${H}$ 四点共面的是\dotfill{(\kh{B})}
    \xx{\includegraphics[width=0.3\textwidth]{bodyxb/src/bo_d4g0jos601uc73fkjrl0_0_291_1650_208_227_0.jpg}}{\includegraphics[width=0.3\textwidth]{bodyxb/src/bo_d4g0jos601uc73fkjrl0_0_583_1665_236_214_0.jpg}}{\includegraphics[width=0.3\textwidth]{bodyxb/src/bo_d4g0jos601uc73fkjrl0_0_847_1665_239_211_0.jpg}}{\includegraphics[width=0.3\textwidth]{bodyxb/src/bo_d4g0jos601uc73fkjrl0_0_1165_1633_213_242_0.jpg}}
    \begin{jieda}
        对于选项$A$，如下图，点$E$、$F$、$H$、$M$确定一个平面，该平面与底面交于$FM$，而点$G$不在平面$EHMF$上，故$E$、$F$、$G$、$H$四点不共面
        \begin{center}                  \includegraphics[width=0.2\linewidth]{bodyxb//src/pingmian(2)B3A.png}
        \end{center}

        对于选项$B$，连结底面对角线$AC$，由中位线定理得$FG//AC$，又$EH//AC$，则$EH//FG$，故$E$、$F$、$G$、$H$四点共面
        \begin{center}                  \includegraphics[width=0.2\linewidth]{bodyxb//src/pingmian(2)B3B.png}
        \end{center}

        对于选项$C$，显然$E$、$F$、$H$所确定的平面为正方体的底面，而点$G$不在该平面内，故$E$、$F$、$G$、$H$四点不共面

        对于选项$D$，如图，取部分棱的中点，顺次连接，得一个正六边形，即点$E$、$G$、$H$确定的平面，该平面与正方体正面的交线为$PQ$，而点$F$不在直线$PQ$上，故$E$、$F$、$G$、$H$四点不共面．
        \begin{center}                  \includegraphics[width=0.2\linewidth]{bodyxb//src/pingmian(2)B3D.png}
        \end{center}
        \end{jieda}
        
    \item 下列说法正确的是\dotfill{(\kh{D})}
    \xx{线段 AB 在平面 $\alpha$ 内,直线 AB 不在 $\alpha$ 内}{平面 $\alpha$ 和 $\beta$ 有时只有一个公共点}{空间三点可以确定一个平面}{若两个平面有不共线的三个公共点, 则两个平面就重合}
    
    \item 在空间四点中，无三点共线是四点共面的\dotfill{(\kh{D})}
    \xx{充分不必要条件}{必要不充分条件}{充分必要条件}{既不充分又不必要条件}
    
    \item 如图，已知 $A\text{ 、 }B\text{ 、 }C$ 在正方体的顶点或棱上，画出过 $A\text{ 、 }B\text{ 、 }C$ 的正方体的截面(保留作图痕迹).
    
    \begin{center}
    \includegraphics[width=0.9\textwidth]{bodyxb/src/bo_d4g0jos601uc73fkjrl0_1_248_784_1146_1046_0.jpg}
    \end{center}

    \begin{jieda}
        \begin{center}
    
            \includegraphics[width=1\linewidth]{asfdbajsfa.png}
        \end{center}
    \end{jieda}
    
    \item 如图， $A\text{ 、 }P$ 是边长为 1 的正方体的顶点， $B\text{ 、 }C$ 在此正方体的棱上，且 $B$ 是所在棱的中点，根据 ${PC}$ 的长度分类讨论过 $A\text{ 、 }B\text{ 、 }C$ 的正方体截面的形状(三角形、四边形、 五边形或者六边形).
    
    \begin{center}
    \includegraphics[width=0.2\textwidth]{bodyxb/src/bo_d4g0jos601uc73fkjrl0_2_209_441_332_313_0.jpg}
    \end{center}

    \begin{jieda}
        ${PC} \in  \left\lbrack  {0,\dfrac{1}{2}}\right\rbrack   \cup   1\}$ : 四边形
        
        ${PC} \in  \left( {\dfrac{1}{2},1}\right)$ : 五边形
    \end{jieda}
    
    \item 如图，已知 $a \cap  \beta  = {EF},A \in  a,{C}$ 、 $B \in  \beta ,{BC}$ 与 ${EF}$ 相交，在图中画出平面 ${ABC}$ 分别与 $\alpha$ 、 $\beta$ 的交线.
    
    \begin{center}
    \includegraphics[width=0.2\textwidth]{bodyxb/src/bo_d4g0jos601uc73fkjrl0_2_242_989_344_217_0.jpg}
    \end{center}
    
    \begin{jieda}
        \begin{center}
            \includegraphics[width=0.25\linewidth]{bodyxb//src/pingmian(2)7.png}
        \end{center}
    \end{jieda}

    \item 已知: 直线 $a\text{ 、 }b\text{ 、 }c$ 相交于 ${P}$ 点,直线 $d$ 与 $a\text{ 、 }b\text{ 、 }c$ 分别相交于 ${A}\text{ 、 }{B}\text{ 、 }{C}$ 三点. 求证: 直线 $a$ 、 $b$ 、 $c$ 与直线 $d$ 共面.
    
    \begin{center}
    \includegraphics[width=0.3\textwidth]{bodyxb/src/bo_d4g0jos601uc73fkjrl0_2_257_1432_381_250_0.jpg}
    \end{center}

    \begin{jieda}
        根据公理2的推论可知$d$和$P$确定一个平面$\alpha$，则$A,B,C,P$均在$\alpha$内，根据公理1可知$a, b, c$均在$\alpha$内.
    \end{jieda}


    \item 如图， $\bigtriangleup  {{ABC}}$ 和 $\bigtriangleup  {{A}}_{1}{{B}}_{1}{{C}}_{1}$ 不在同一个平面内，如果三直线 ${A}{{A}}_{1}\text{ 、 }{B}{{B}}_{1}\text{ 、 }{{{CC}}}_{1}$ 两两相交，求证: 三直线 $A{A}_{1}\text{ 、 }B{B}_{1}\text{ 、 }C{C}_{1}$ 交于一点.

    \begin{center}
    \includegraphics[width=0.3\textwidth]{bodyxb/src/bo_d4g0jos601uc73fkjrl0_3_246_337_358_238_0.jpg}
    \end{center}

    \begin{jieda}
        设 $B{B}_{1}$ 与 $C{C}_{1},C{C}_{1}$ 与 $A{A}_{1},A{A}_{1}$ 与 $B{B}_{1}$ 分别确定平面 $\alpha ,\beta ,\gamma ,A{A}_{1}$ 与 $B{B}_{1}$ 的交点为 $P$
        
        因为 $P \in  A{A}_{1},P \in  B{B}_{1}$，$A{A}_{1} \subset  \beta ,B{B}_{1} \subset  \alpha$ ,所以 $P \in  \alpha ,P \in  \beta$ ,即 $P \in  \alpha  \cap  \beta$ .
        
        又 $\alpha  \cap  \beta  = C{C}_{1}$ ,所以 $P \in  C{C}_{1}$ .所以三条直线 $A{A}_{1},B{B}_{1},C{C}_{1}$ 交于一点
    \end{jieda}

\end{enumerate}

C 组:
\begin{enumerate}
    \item 到空间不共面的 4 个点距离都相等的平面有\dotfill{(\kh{A})}
    \xx{7 个}{6 个}{4 个}{3 个}
    \begin{jieda}
        将空间不共面的四个点看作四面体的顶点,
        
        第一类:平行于底面，且过测棱的中点，此时由 4 种情况符合要求，如图:
        
        \begin{center}
        \includegraphics[width=0.2\textwidth]{bodyxb/src/bo_d4hptm3ef24c739bif30_0_243_609_265_301_0.jpg}
        \end{center}
        
        第二类:一条棱的中点于底面的一条中位线所成的平面，此时有 3 种情况符合要求，如图
        
        \begin{center}
        \includegraphics[width=0.2\textwidth]{bodyxb/src/bo_d4hptm3ef24c739bif30_0_243_978_265_302_0.jpg}
        \end{center}
        
        故总共有 7 种情况符合要求.
    \end{jieda}
    
    \item 空间有 6 个点，其中任意三点不共线，且有五个点共面，则这 6 个点最多可以确定\blkda{11} 个平面.

    \begin{jieda}
        记 5 个共面的点分别为 $a,b,c,d,e$ ，不共面的第 6 个点为 $O$ ，则从 $a,b,c,d,e$ 中随机选取两个点，连同第 6 个点，即可构成一个平面，此时共有
        
        $\left( {abO}\right) ,\left( {acO}\right) ,\left( {adO}\right) ,\left( {aeO}\right) ,\left( {bcO}\right) ,\left( {bdO}\right) ,\left( {beO}\right) ,\left( {cdO}\right) ,\left( {ceO}\right)$ 共有 10 个平面，(abcde) 这 5 个点确定一个平面，故最多可以确定 10+1=11 个平面
    \end{jieda}
    
    \item 金字塔在埃及和美洲等地均有分布，现在的尼罗河下游，散布着约 80 座金字塔遗迹，大小不一, 其中最高大的是胡夫金字塔, 如图, 胡夫金字塔可以近似看做一个正四棱锥, 则该正四棱锥的 5 个面所在的平面将空间分成\blkda{23}部分 (用数字作答).
    
    \begin{center}
    \includegraphics[width=0.3\textwidth]{bodyxb/src/bo_d4g0jos601uc73fkjrl0_3_244_1253_400_272_0.jpg}
    \end{center}
    \begin{jieda}
        假想一个没有上顶的正方体，该正方体会把空间分割成 18 块，把四面进行极限倾斜相交，如图所示，
        
        \begin{center}
        \includegraphics[width=0.2\textwidth]{bodyxb/src/bo_d4hptm3ef24c739bif30_1_243_889_281_260_0.jpg}
        \end{center}
        
        在倾斜的过程中，在不管底面的情况下，4 个侧面在顶点以下的“水平范围”内最多可以切割出 9 个空间，与没有倾斜极限的情况一样，多出来的空间是交叉的切割出来的空间，
        
        在空间上是对称的，四个倾斜的侧面在空间中的延伸还是这样的倾斜侧面，
        
        如图所示的对称的锥面同样会切割出 $g$ 个空间，即顶点之上的 4 个延伸的倾斜的面同样会切割出 9 个空间.
        
        但是四个空间和下面的四个倾斜的侧面切出的是同一个，即标记“ $\times$ ”的位置，
        
        \begin{center}
        \includegraphics[width=0.3\textwidth]{bodyxb/src/bo_d4hptm3ef24c739bif30_1_239_1681_359_281_0.jpg}
        \end{center}
        
        所以在 18 的基础上加 9 减 4 ，即结果是 23.
    \end{jieda}

\end{enumerate}

\clearpage
\subsection*{平面及其基本性质 (3)}
\begin{enumerate}
    \item 两个相等的角在直观图中\blkda{不一定}相等；两条相等的线段在直观图中\blkda{不一定} 相等. (填 “一定” “不一定”)

    \item 如果正三角形 ${ABC}$ 的边长为 $a$ ，那么 $\bigtriangleup  {ABC}$ 的水平放置的直观图 $\bigtriangleup  {A}^{\prime }{B}^{\prime }{C}^{\prime }$ 的面积为\blkda{$\dfrac{\sqrt{6}}{16}a^2$}，已知 $\bigtriangleup  {ABC}$ 的平面直观图 $\bigtriangleup  {A}^{\prime }{B}^{\prime }C$ 是边长为 $a$ 的正三角形，那么原$\bigtriangleup  {ABC}$ 的面积为\blkda{$\dfrac{\sqrt{6}}{2}a^2$}.

    \item 若右图表示水平放置的 $\bigtriangleup  {ABC}$ 的直观图(其中 $D$ 为 ${BC}$ 的中点)，则在原来的 ${\Delta ABC}$ 中，四条线段中最长的是\blkda{AC}.
    
    \begin{center}
    \includegraphics[width=0.2\textwidth]{bodyxb/src/bo_d4g0jtbef24c73ek1ku0_0_393_724_276_204_0.jpg}
    \end{center}
    
    \item 一梯形的直观图是一个如图所示的等腰梯形，且该梯形面积为 $\sqrt{2}$ ，则原梯形的面积为\blkda{4}.
    
    \begin{center}
    \includegraphics[width=0.2\textwidth]{bodyxb/src/bo_d4g0jtbef24c73ek1ku0_0_996_770_245_152_0.jpg}
    \end{center}

    \item 利用斜二测画法得到的:
    \begin{dingautolist}{172}
        \item 三角形的直观图一定是三角形； \item 正方形的直观图一定是菱形；
        \item 等腰梯形的直观图可以是平行四边形； \item 菱形的直观图一定是菱形.
    \end{dingautolist}
    以上结论正确的是\blkda{\ding{172}}.
    
    \item 如图，模块 \ding{172}-\ding{176} 均由 4 个棱长为 1 的小正方体构成，模块 \ding{177} 由 15 个棱长为 1 的小正方体构成. 现从模块 \ding{172}-\ding{176} 中选出三个放到模块 \ding{177} 上，使得模块 \ding{177} 成为一个棱长为 3 的大正方体. 则十种可能的选择方案中，能够完成任务的为\blkda[4]{\ding{172}\ding{173}\ding{176}或\ding{172}\ding{175}\ding{176}或\ding{173}\ding{174}\ding{175}}.
    
    \begin{center}
    \includegraphics[width=0.8\textwidth]{bodyxb/src/bo_d4g0jtbef24c73ek1ku0_0_241_1690_981_223_0.jpg}
    \end{center}

    \item 若长方体的全面积为 22，棱长之和为 24，则其体对角线的长为\dotfill{(\kh{D})}
    \xx{$\sqrt{11}$}{$\sqrt{12}$}{$\sqrt{13}$}{$\sqrt{14}$}
    
    \item 在正方体 ${ABCD} - {A}_{1}{B}_{1}{C}_{1}{D}_{1}$ 中，点 $P$ 在侧面 ${BC}{C}_{1}{B}_{1}$ 及其边界上运动，并且总是保持 ${AP} \bot  {B{D}_{1}}$ ，则动点 $P$ 的轨迹是\dotfill{(\kh{B})}
    \xx{一个点}{线段}{正三角形}{圆}
    
    \item 在从长方体同一顶点出发的三条棱上各取一点，过它们作截面，所得截面的形状一定是\dotfill{(\kh{C})}
    \xx{等边三角形}{直角三角形}{锐角三角形}{钝角三角形}
    
    \item 在水平放置的平面 $a$ 上有一个四边形,当用斜二测画法画出它的直观图时, 它被画成边长为 $2{{cm}}$ 的正方形，并且底边在水平位置，如图所示. 画出这个四边形的真实形状，并求它的面积.
    
    \begin{center}
    \includegraphics[width=0.4\textwidth]{bodyxb/src/bo_d4g0jtbef24c73ek1ku0_1_872_990_521_194_0.jpg}
    \end{center}

    \begin{jieda}
        $S = 8\sqrt{2} cm^2$
    \end{jieda}
    
    \item 如图，已知四边形 ${OABC}$ ，画出它的直观图.
    
    \begin{center}
    \includegraphics[width=0.2\textwidth]{bodyxb/src/bo_d4g0jtbef24c73ek1ku0_1_209_1338_343_337_0.jpg}
    \end{center}

    \item 在水平放置的平面 $a$ 上，画一个相邻两边长分别为 $2{{cm}}$ 、 $3{{cm}}$ ，且它们的夹角为 ${60}^{ \circ  }$ 的平行四边形的直观图.

    \item 右图是某个水平放置的平面图形的直观图，画出原来的平面图形.
    
    \begin{center}
    \includegraphics[width=0.3\textwidth]{bodyxb/src/bo_d4g0jtbef24c73ek1ku0_2_948_325_459_269_0.jpg}
    \end{center}
\end{enumerate}


\section*{直线与直线的位置关系}
\subsection*{直线与直线的位置关系（1）}

A 组：
\begin{enumerate}

    \item 如果一个角的两边和另一个角的两边分别平行, 那么这两个角的
关系是\blkda[3]{相等或互补}.
    
    
    \item 已知空间四边形 $ {ABCD} $ (即 $ {A}\text{、}{\;B}\text{、}{C}\text{、}{D} $ 四点不共面) 各边中点
分别为 $ E $ 、 $ F $ 、 $ G $ 、 $ H $ ,则四边形 $ {EFGH} $ 是\blkda[2]{平行四边}形.

\begin{center}
    \centering
    \includegraphics[width=0.25\linewidth]{imagiasfygiuahgfe.png}
\end{center}

    \item 判断下列命题的真假,真的打“√”,假的打“×”：
    \begin{enumerate}
        \item 平行于同一直线的两条直线平行 \blkda{\ding{51}}
        \item 过直线外一点,有且只有一条直线与已知直线平行 \blkda{\ding{51}}
        \item 与已知直线平行且距离等于定长的直线只有两条 \blkda{\ding{55}}
        \item 若一个角的两边分别与另一个角的两边平行,那么这两个角相等\blkda{\ding{55}}
        \item 若两条相交直线和另两条相交直线分别平行,那么这两组直线所成的锐角(或直角)相等 \blkda{\ding{51}}

    \end{enumerate}

\end{enumerate}

B组：
\begin{enumerate}
    \item 若 $ {AB}//{A}^{\prime }{B}^{\prime },{BC}//{B}^{\prime }{C}^{\prime } $ ,且 $ \angle {ABC} = {45}^{ \circ  } $ ,则 $ \angle {A}^{\prime }{B}^{\prime }{C}^{\prime } $ 等于\blkda[2]{$45^{\circ}$或$135^{\circ}$}.
    
    \item 下列命题中,正确的是\blkda{(2)(4)}. (在横线上填写所有正确说法的编号)
    \begin{enumerate}
        \item 有三个角是直角的四边形是矩形；
        \item 有四个角是直角的四边形是矩形；
        \item 四边相等且四个角也相等的四边形一定是正方形；
        \item 两条对角线互相平分的四边形是平行四边形.
    \end{enumerate}
    

    \item 已知空间中两个平行四边形 $ {ABCD} $ 和 $ {EFGH} $ ,若 $ {AB}\overset{//}{ = }{EF},{BC}\overset{//}{ = }{FG} $ ,则下列说法中一定正确的是\blkda{(1)(4)}.
    \begin{enumerate}
        \item  $ {S}_{ABCD} = {S}_{EFGH} $ ；
        \item $ {AC} = {EG} $ ；
        \item 直线 $ {BD} $ 与直线 $ {FH} $ 可以确定一个平面；
        \item 直线 $ {CD} $ 和直线 $ {GH} $ 可以确定一个平面.
    \end{enumerate}

    \item 若空间中四条直线 $ {l}_{1} $ 、 $ {l}_{2} $ 、 $ {l}_{3} $ 、 $ {l}_{4} $ 满足 $ {l}_{1}\bot {l}_{2} $ , $ {l}_{2}//{l}_{3} $ , $ {l}_{3}\bot {l}_{4} $ ,则下列结论一定正确的是\blkda{4,6}.\\
    (1) $ {l}_{1}//{l}_{3} $ 、(2) $ {l}_{1}//{l}_{4} $ 、(3) $ {l}_{2}//{l}_{4} $ 、(4) $ {l}_{1}\bot {l}_{3} $ 、(5) $ {l}_{1}\bot {l}_{4} $ 、(6) $ {l}_{2}\bot {l}_{4} $ .
    
    \item 如图所示,在空间四边形 $ {{ABCD}} $ 中,E、F、G、H 分别是 $ {{AB}} $ 、 $ {{BC}} $ 、 $ {{CD}} $ 、 $ {{DA}} $的中点.
    \begin{enumerate}
        \item 求证：EFGH 是平行四边形；
        \begin{center}
    
            \includegraphics[width=0.25\linewidth]{fhyauhfa'uoeffa.png}
        \end{center}
        \begin{jieda}
            因为$ EH \overset{/\!/}{=\!\!=} \dfrac{1}{2}BD,FG \overset{/\!/}{=\!\!=} \dfrac{1}{2}BD$,所以$EH \overset{/\!/}{=\!\!=} FG$,
            故为平行四边形.
        \end{jieda}

        \item 若 $ {EFGH} $ 为菱形,则 $ {AC} $ 、 $ {BD} $ 满足的条件为\blkda{$AC=BD$}
        
        \item 若 $ {{EFGH}} $ 为矩形,则 $ {AC} $ 、 $ {BD} $ 满足的条件为\blkda{$AC\perp BD$}
        
        \item 若 $ {EFGH} $ 为正方形,则 $ {AC} $ 、 $ {BD} $ 满足条件为\blkda[4]{$AC\perp BD$且$AC=BD$}
        
        \item 若 $ {{AC}} + {{BD}} = a $ , $ {{AC}} \cdot  {{BD}} = b $ ,则 $ E{G}^{2} + H{F}^{2} = $ \blkda{$\dfrac{1}{2}a^2-b$}
    \end{enumerate}
    
    \item 如图,在空间四边形 $ {{ABCD}} $ 中,点 $ {E} $ 、 $ {H} $ 分别是边 $ {{AB}} $ 、$ {{AD}} $ 上的点, $ {F} $ 、 $ {G} $ 分别是边 $ {{BC}} $ 、 $ {{CD}} $ 上的点. 若点 $ {E} $ 、 $ {H} $ 分别是边 $ {{AB}} $ 、 $ {{AD}} $ 的中点,且 $ \dfrac{CF}{CB} = \dfrac{CG}{CD} = \dfrac{2}{3} $ ,则
        \dotfill{(\kh{D})}
        \begin{center}
    
            \includegraphics[width=0.25\linewidth]{imaafha'oisfge.png}
        \end{center}
    \xx{EF与 $ {GH} $ 互相平行;}
    {EF与GH 异面;}
    {EF与 $ {GH} $ 的交点 $ {M} $ 可能在直线 $ {{AC}} $ 上,也可能不在直线 $ {{AC}} $ 上；}
    {EF与 $ {GH} $ 的交点 $ M $ 一定在直线 $ {AC} $ 上.}
    \begin{jieda}
        因为$EF\subset 平面ABC，HG\subset 平面ACD$，由公理3可知交点$M$一定在两个面的交线$AC$上.
    \end{jieda}

   
    \item 下列命题中,正确结论有\blkda{2,3}.
    \begin{enumerate}
        \item 如果一个角的两边与另一个角的两边分别平行,那么这两个角相等；
        \item 如果两条相交直线和另两条相交直线分别平行,那么这两组直线所成的锐角 (或直角) 相等;
        \item 如果两条直线同平行于第三条直线,那么这两条直线互相平行.     
    \end{enumerate}

    \item 如果空间四边形的对角线相等,那么顺次联结它的各边中点所成的四边形一定是\dotfill{(\kh{C})}
    \xx{梯形}{长方形}{菱形}{正方形}

    
    \item 在空间四边形 $ {ABCD} $ 中, $ P $ 、 $ Q $ 分别是边 $ {AB} $ 、 $ {CD} $ 的中点,设 $ {BD} + {AC} = {2a} $ ,则\dotfill{(\kh{C})}
    \xx{$ {PQ} > a $}
    {$ {PQ} = a $}
    {$ {PQ} < a $}
    {$ {PQ} $ 与 $ a $ 的大小不确定} 
    
    \item 在空间四边形 $ {{ABCD}} $ 的边 $ {{AB}} $ 、 $ {{BC}} $ 、 $ {{CD}} $ 、 $ {{DA}} $ 上分别取 $ {E} $ 、 $ {F} $ 、 $ {G} $ 、 $ {H} $ 四点,如果 $ {{EF}} $ 与 $ {{HG}} $ 交于点 $ {M} $ ,则 \dotfill{(\kh{A})}
    \xx{M一定在直线 AC 上}
    {M一定在直线 $ {{BD}} $ 上}
    {$ {M} $ 可能在 $ {{AC}} $ 上,也可能在 $ {{BD}} $ 上}
    {$ {M} $ 不在 $ {{AC}} $ 上,也不在 $ {{BD}} $ 上} 

    \item 如图,在长方体 $ {ABCD} - {A}_{1}{B}_{1}{C}_{1}{D}_{1} $ 中, $ M $ 为 $ {B}_{1}B $ 上一点,试过点 $ M $ 作平面 $ {D}_{1}{DC}{C}_{1} $的对角线 $ D{C}_{1} $ 的平行线,并说明理由.

    \begin{center}

        \includegraphics[width=0.25\linewidth]{fah'lushdf'afas.png}
    \end{center}

    \begin{jieda}
        在$AB$上找作点$N$，使得$MN\parallel AB_1$，根据直线平行的传递性可知$MN\parallel DC_1$
    \end{jieda}


    \item 如图, $ \bigtriangleup {{ABC}} $ 和 $ \bigtriangleup {{A}}_{1}{{\;B}}_{1}{{C}}_{1} $ 的对应顶点连线, $ {{{AA}}}_{1} $ 、 $ {{{BB}}}_{1} $ 、$ {{{CC}}}_{1} $ 交于点 $ {O},{O} $ 在平面 $ {{ABC}} $ 和平面 $ {{A}}_{1}{{\;B}}_{1}{{C}}_{1} $ 之间,且$ \dfrac{{{AO}}}{{{A}}^{\prime }{O}} = \dfrac{{{BO}}}{{{B}}^{\prime }{O}} = \dfrac{{{CO}}}{{{C}}^{\prime }{O}} = \dfrac{2}{3}. $
    \begin{enumerate}
        \item 求证： $ \angle {CAB} = \angle {C}_{1}{A}_{1}{B}_{1} $ ；
        \item 求: $ {S}_{\bigtriangleup {ABC}} : {S}_{\bigtriangleup {A}_{1}{B}_{1}{C}_{1}} $ 的值.
    \end{enumerate}
    \begin{center}

        \includegraphics[width=0.25\linewidth]{f'ahsufa'fad.png}
    \end{center}
\end{enumerate}

    \begin{jieda}
        \begin{enumerate}
            \item 因为$\triangle ACO \sim \triangle A_1C_1O$，所以$AC\parallel A_1C_1$,因为$\triangle ABO \sim A_1B_1O$，所以$AB\parallel A_1B_1$.由等角定理可得，$\angle{CAB}=\angle{C_1A_1B_1}$

            \item $\triangle ABC \sim \triangle A_1B_1C_1$,因此面积之比是边长之比的平方，故4:9.
            
        \end{enumerate}
    \end{jieda}
\subsection*{直线与直线的位置关系（2）}

A 组：
\begin{enumerate}
    \item 分别在两个平面内的两条直线的位置关系是\blkda[4]{异面、相交、平行}；不平行的两条直线的位置关系是\blkda[3]{相交或异面}；没有公共点的两条直线的位置关系\blkda[3]{平行或异面}.

    \item 若两异面直线所成的角为 $ \theta $ ,则角 $ \theta $ 的范围是\blkda[3]{$(0,\dfrac{\pi}{2}]$}.
    
    \item 两条异面直线是指
    \dotfill{(\kh{D})}
    \xx{在空间不相交的两条直线}
    {分别位于两个不同平面内的两条直线}
    {某平面内的一条直线和这个平面外的一条直线}
    {不在同一平面内的两条直线} 
    
    \item 分别在两个相交平面内的两条直线的位置关系是
    \dotfill{(\kh{D})}
    \xx{异面}{平行}{相交}{异面、平行或相交}
    
    \item 分别在两个平行平面内的两条直线的位置关系是
    \dotfill{(\kh{D})}
    \xx{异面}{平行}{相交}{异面或平行}
    
    \item 分别与两条异面直线都相交的两条直线的位置关系是
    \dotfill{(\kh{B})}
    \xx{平行或相交}{相交或异面}{平行或异面}{均有可能}    

    \item 如下图,正方体 $ {ABCD} - {A}_{1}{B}_{1}{C}_{1}{D}_{1} $ 中, $M\text{、}N\text{、}P\text{、}Q\text{、}R $ 、$ S $ 分别是 $ {AB}\text{、}{BC}\text{、}{C}_{1}{D}_{1}\text{、}{C}_{1}C\text{、}{A}_{1}{B}_{1}\text{、}{B}_{1}B $ 的中点,则下列判断：(1) $ {PQ} $ 与 $ {RS} $ 共面；(2) $ {MN} $ 与 $ {RS} $ 共面；(3) $ {PQ} $与 $ {MN} $ 共面；则正确的结论是\blkda{(1)(3)}.
    
    \begin{center}

        \includegraphics[width=0.25\linewidth]{i;ugyf;iyffsa'i.png}
    \end{center} 
    
\end{enumerate}


B组：
\begin{enumerate}
    \item 已知 $ a $ 、 $ b $ 是异面直线,直线 $ c//a $ ,那么直线 $ c $ 与 $ b $ 
    \dotfill{(\kh{D})}
    \xx{一定是异面直线}{一定是相交直线}{不可能是相交直线}{不可能是平行直线}

    \item 已知 $ {a} $ 、 $ {b} $ 、 $ {c} $ 是空间三条直线,有下列四个命题,其中正确命题的个数是\blkda{1个}.
    \begin{enumerate}
        \item 若 $ a\bot b,b\bot c $ ,则 $ a\bot c $
        \item 若 $ a $ 、 $ b $ 异面, $ b $ 、 $ c $ 异面,则 $ a $ 、 $ c $ 异面
        \item 若 $ a $ 、 $ b $ 共面, $ b $ 、 $ c $ 共面,则 $ a $ 、 $ c $ 共面
        \item 若 $ a $ 、 $ b $ 平行, $ b $ 、 $ c $ 平行,则 $ a $ 、 $ c $ 平行
    \end{enumerate}

    \item 正方体 $ {{ABCD}} - {{A}}_{1}{{\;B}}_{1}{{C}}_{1}{{D}}_{1} $ 中, $ {{A}}_{1}{\;A} $ 和 $ {{D}}_{1}{{C}}_{1} $ 所成的角是\blkda{ $\dfrac{\pi}{2}$}, $ {{AC}} $ 和$ {{A}}_{1}{{B}}_{1} $ 所成的角是\blkda{ $\dfrac{\pi}{4}$}, $ {{{DA}}}_{1} $ 和 $ {{AC}} $ 所成的角是\blkda{ $\dfrac{\pi}{3}$}, $ {{AC}} $ 与 $ {{D}}_{1}{{B}}_{1} $所成的角是\blkda{ $\dfrac{\pi}{2}$}, $ {{{AD}}}_{1} $ 和 $ {{{DC}}}_{1} $ 所成的角是\blkda{$\dfrac{\pi}{3}$ }.
    
    \item 在长方体 $ {ABCD} - {A}_{1}{B}_{1}{C}_{1}{D}_{1} $ 中, $ {AB} = 5,{BC} = 4,{C{C}_{1}} = 3 $ ,则 $ {B}_{1}C $ 与 $ D{D}_{1} $ 所成的角的大小为\blkda{$arctan \dfrac{4}{3}$}, $ {BC} $ 与 $ {A}_{1}{C}_{1} $ 所成角的大小为\blkda{$arctan\dfrac{5}{4}$}, $ {B}_{1}C $ 与 $ A{D}_{1} $所成角的大小为\blkda{$arccos\dfrac{7}{25}$}

    \item 右下图是一个正方体的展开图,在原正方体中,有下列命题：
    \begin{enumerate}
        \item AB 与 CD 所在直线垂直；
        \item CD 与 EF 所在直线平行;
        \item $ {{AB}} $ 与 $ {{MN}} $ 所在直线成 $ {60}^{ \circ  } $ 角；
        \item MN 与 EF 所在直线异面
    \end{enumerate}
    其中正确命题的序号是\blkda{(3)(4)}.
    \begin{center}

        \includegraphics[width=0.25\linewidth]{imasdufha;iuge.png}
    \end{center}

    \item 在正方体 $ {ABCD} - {A}_{1}{B}_{1}{C}_{1}{D}_{1} $ 中,与 $ A{D}_{1} $ 成 $ {60}^{ \circ  } $ 角的面对角线共有\blkda{8}条.
    
    \item 已知正方体 $ {{ABCD}} - {{A}}_{1}{{B}}_{1}{{C}}_{1}{{D}}_{1} $ 中, $ {M} $ 、 $ {N} $ 分别是 $ {{{BB}}}_{1} $ 、 $ {{{CC}}}_{1} $ 的中点,则异面直线AM 与 BN 所成的角的大小是\blkda{$arccos\dfrac{1}{5}$}.

    \item 在棱长为 1 的正方体 $ {{ABCD}} - {{A}}_{1}{{B}}_{1}{{C}}_{1}{{D}}_{1} $ 中, $ {O} $ 是正方形 $ {{ABCD}} $ 的中心, $ {E} $ , $ {F} $分别是 $ {{AB}},{{BC}} $ 中点.
    \begin{enumerate}
        \item 异面直线 $ {{A}}_{1}{{D}}_{1} $ 和 $ {{CD}} $ 的距离是\blkda{1};
        
        \item $ {{A}}_{1}{{B}}_{1} $ 与 $ {{{BC}}}_{1} $ 的距离是\blkda{$\dfrac{\sqrt{2}}{2}$}；

        \item 异面直线 $ {{C}}_{1}{O} $ 和 $ {{EF}} $ 的距离是\blkda{$\dfrac{\sqrt{2}}{4}$}.
        
    \end{enumerate}

    \item 如图,在空间四边形 $ {{ABCD}} $ 中, $ {{AC}}\bot {{BD}} $ , $ {{AC}} = {{BD}} $ , 若 $ {E},{F} $ 分别是 $ {{AB}},{{CD}} $ 的中点,则 $ {{EF}} $ 与 $ {{AC}} $ 所成的角的大小为\blkda{$\dfrac{\pi}{4}$}.

    \item 若 $ {P} $ 为异面直线 $ {a} $ 、 $ {\;b} $ 外一点,则过 $ {P} $ 与 $ {a} $ 、 $ {\;b} $ 同时：
    \begin{enumerate}
        \item 平行的直线有\blkda{0}条；
        \item 垂直的直线有\blkda{1}条；
        \item 相交的直线有\blkda{0或1}条.
    \end{enumerate}
    
    \item 异面直线 $ {a} $ 、 $ {b} $ 分别在平面 $ {a} $ 、 $ {\beta } $ 内, $ {a} \cap  {\beta } = {L} $ ,则直线 $ {L} $ 与 $ {a} $ 、 $ {b} $ 的位置关系是\dotfill{(\kh{B})}
    \xx{与a、b都相交} {至少与a、b中的一条相交}{与 a、b 都不相交}{至多与 $ {a}\text{、}{\;b} $ 中的一条相交}

    \item 如图所示,在空间四边形 $ {{ABCD}} $ 中,E、F、G、H 分别是 $ {{AB}} $ 、 $ {{BC}} $ 、 $ {{CD}} $ 、 $ {{DA}} $的中点. 
    \begin{enumerate}
        \item 若 $ {{AC}} $ 与 $ {{BD}} $ 成 $ {60}^{ \circ  } $ 角,且 $ {{AC}} = {{BD}} = a $ ,则 $ {{EFGH}} $ 的面积为\blkda{$\dfrac{\sqrt{3}}{8}a^2$}
        \item 若 $ M $ 、 $ N $ 分别是对角线 $ {AC} $ 、 $ {BD} $ 的中点. 求证： $ {EG} $ 、 $ {FH} $ 、$ {MN} $ 交于一点且在该点互相平分.
    \end{enumerate}
    \begin{center}

        \includegraphics[width=0.7\linewidth]{f;kshdfuhsdf.png}
    \end{center}

        \begin{jieda}
            在平行四边形EHGF中，EG与FH是平行四边形的对角线，因此互相平分；在平行四边形ENGM中，EG与NM为平行四边形的对角线，因此互相平分.
        \end{jieda}

    \item 如图,在空间四边形 $ {{ABCD}} $ 中,点 $ {E} $ 、 $ {H} $ 分别是边 $ {{AB}} $ 、 $ {{AD}} $ 上的点, $ {F} $ 、 $ {G} $分别是边 $ {{BC}}\text{、}{{CD}} $ 上的点. 若 $ {{EFGH}} $ 是平行四边形,且 $ {{BD}} = a,{{AC}} = b,{{AC}} $ 与 $ {{BD}} $所成的角为 $ \theta $ ,四边形 EFGH 面积的最大值为\blkda{$\dfrac{1}{4}absin\theta$}.
    \begin{jieda}
        因为EFGH是平行四边形，因此$AE:EB=$
    \end{jieda}

    \item 如图,空间四边形 $ {ABCD} $ 中, $ {AD} = {BC} = 4,E\text{、}F $ 分别是 $ {AB}\text{、}{CD} $ 的中点, $ {EF} = 2\sqrt{3} $ .则 AD、BC 所成角为\blkda{$\dfrac{\pi}{3}$}.

    \item 如图所示,在正方体 $ {ABCD} - {A}_{1}{B}_{1}{C}_{1}{D}_{1} $ 中, $ O $ 是底面 $ {ABCD} $的中心, $ E $ 是 $ C{C}_{1} $ 的中点. 那么异面直线 $ {EO} $ 和 $ {D}_{1}A $ 所成的角的余弦值等于\blkda{$\dfrac{\sqrt{6}}{3}$}.
    
    \begin{center}

        \includegraphics[width=0.25\linewidth]{imsfuh'afueafage.png}
    \end{center}

    \item 对于已知直线 $ a $ ,假设直线 $ b $ 同时满足下列三个条件：(1)与 $ a $ 成异面直线；(2)与 $ a $ 的夹角为定值 $ \theta $ ；(3)与 $ a $ 的距离为定值 $ d $ .那么,这样的直线 $ b $ 有\blkda{无数}条.

    
    \item 在正方体的 12 条棱、12 条面对角线中,总共可以组成\blkda{126}对异面直线.
    \begin{jieda}
        棱与棱可以形成$\dfrac{12*4}{2}=24$对异面直线，线与线可以形成$\dfrac{12*5}{2}=30$对异面直线，棱与线可以形成$6*12=72$对异面直线，
    \end{jieda}
    
    \item 若 $ a,b $ 为异面直线, $ m \cap  a = A,m \cap  b = B,n \cap  a = M,n \cap  b = N $ ,求证： $ m,n $是异面直线.
    \begin{jieda}
        因为直线$a$与直线$b$异面，所以$A\notin b$，从而A与b可确定平面$\alpha$.因为$B\in b, b\subset\alpha$，所以$B\in \alpha$，又因为$A\in \alpha$，所以$AB \subset \alpha$.因为直线$a$与直线$b$异面，所以$M\notin \alpha$.因为$N\in b\subset \alpha$,且$N\notin m$，由异面直线的判定定理，n与m异面.
        \begin{center}
            \includegraphics[width=0.25\linewidth]{a'huf'auhfdf.png}
        \end{center}
    \end{jieda}
    
    \item 如图,在正方体 $ {{ABCD}} - {{A}}_{1}{{\;B}}_{1}{{C}}_{1}{{D}}_{1} $ 中, $ {E} $ 是 $ {{B}}_{1}{{C}}_{1} $ 的中点, $ {O} $ 是正方形 $ {{A}}_{1}{{B}}_{1}{{C}}_{1}{{D}}_{1}{R} $的中心,连接 $ {{AO}},{{CE}} $ ,求异面直线 $ {{AO}} $ 与 $ {{CE}} $ 所成的角的大小.

    \begin{center}

        \includegraphics[width=0.25\linewidth]{fhsuhfajsdf.png}
    \end{center}

    \begin{jieda}
        取$CD$的中点M，可证得$\angle{AOM}$即为异面直线所成角，$|AO|=\dfrac{\sqrt{6}}{2},|MO|=\dfrac{\sqrt{5}}{2},|AM|=\dfrac{\sqrt{5}}{2}$，由余弦定理可得$cos\angle{AOM=\dfrac{\sqrt{30}}{10}}$，综上异面直线AO与CE所成角为$arccos \dfrac{\sqrt{30}}{10}$.
        \begin{center}
    
            \includegraphics[width=0.25\linewidth]{asoufha'oiusf.png}
        \end{center}
    \end{jieda}

    
    \item 在正方体 $ {ABCD} - {A}_{1}{B}_{1}{C}_{1}{D}_{1} $ 中,正方形 $ {AB}{B}_{1}{A}_{1} $ 和正方形 $ {BC}{C}_{1}{B}_{1} $ 的中心分别为 $ O $和 $ {O}_{1} $ ,试判断直线 $ O{O}_{1} $ 与直线 $ B{D}_{1} $ 的位置关系,并证明你的结论.
    \begin{jieda}
        异面.记平面$A_1BC_1$为$\alpha$，因为$OO_1\subset\alpha, B\in \alpha 且 B\notin OO_1，D\notin \alpha$，由异面直线判定定理可知，两直线异面.
        \begin{center}
    
            \includegraphics[width=0.25\linewidth]{imauohf'aouhfage.png}
        \end{center}
    
    \end{jieda}

    \item 已知正方体 $ {ABCD} - {A}_{1}{B}_{1}{C}_{1}{D}_{1} $ 的棱长为 $ a,M\text{、}N $ 分别为 $ {A}_{1}{B}_{1}\text{、}{B}_{1}B $ 的中点,求 $ {AM} $ 与 $ {CN} $ 所成的角.

    \begin{center}

        \includegraphics[width=0.25\linewidth]{imfsdhfffage.png}
    \end{center}

    \begin{jieda}
        取$D_1C_1$的中点$M_1$，$AA_1$的中点$N_1$，可证明$\angle{N_1DM_1}$即为异面直线的夹角，利用余弦定理可得$arccos \dfrac{2}{5}$
        \begin{center}
    
            \includegraphics[width=0.25\linewidth]{ahfa'ouhfaaf.png}
        \end{center}
    \end{jieda}
\end{enumerate}



\section*{直线与平面的位置关系}
\subsection*{直线与平面的位置关系（1）}

知识点：线面平行

A组：
\begin{enumerate}
    \item 判断下列命题是否正确:
        \begin{enumerate}
            \item 过平面外一点可作无数条直线与这个平面平行. \blkda{\ding{51}}
            
            \item 若直线 $ l ⊄ a $ ,则 $ l $ 不可能与 $ \alpha $ 内无数条直线相交. \blkda{\ding{55}}
            
            \item 若直线 $ l $ 与平面 $ \alpha $ 不平行,则 $ l $ 与 $ \alpha $ 内任一直线都不平行.\blkda{\ding{55}}
            
            \item 经过两条平行线中一条直线的平面平行于另一条直线.\blkda{\ding{55}}
            
            \item 若平面 $ \alpha $ 内有一条直线和直线 $ l $ 异面,则 $ l ⊄ a $ . \blkda{\ding{51}}
        \end{enumerate}
 
    \item 已知平面 $ \alpha $ 与直线 $ l $ ,则 “ $ l//\alpha $ ” 是 “直线 $ l $ 与平面 $ \alpha $ 无公共点” 的\dotfill{(\kh{C})}
    \xx{充分不必要条件} {必要不充分条件}{充要条件}{既不充分也不必要条件}

    \item 下列命题中,正确的是\dotfill{(\kh{D})}
    \xx{若直线 $ a $ 与平面 $ \alpha $ 平行,则 $ a $ 平行于 $ \alpha $ 内的任何直线}
    {若两直线 $ a,b $ 都与平面 $ {\alpha } $ 平行,则 $ {a}//{b} $}
    {若直线 $ a $ 平行于平面 $ \alpha $ ,直线 $ b $ 在平面 $ \alpha $ 内,则 $ a//b $}
    {若直线 $ l $ 与平面 $ \alpha $ 平行,则平面 $ \alpha $ 内有无数条直线与 $ l $ 平行}
    
    \item 如图,在下列四个正方体中, $ {A},{B} $ 为正方体的两个顶点, $ {M},{N},{Q} $ 为所在棱的中点, 则在这四个正方体中, 直线 $ {{AB}} $ 与平面 $ {{MNQ}} $ 平行的有\blkda{3}个.
    \begin{jieda}
        (1)中令$NM$的中点为T,则$AB$与$QT$相交（因为共面且不平行）;(2)中$AM // MQ$；(3)中$AB // MQ$;（4）中$AB//NQ$
    \end{jieda}

    \begin{center}

        \includegraphics[width=0.75\linewidth]{imaarfharhfafge.png}
    \end{center}

\end{enumerate}

B组：
\begin{enumerate}
    \item 在正方体 $ {{ABCD}} - {{A}}_{1}{{B}}_{1}{{C}}_{1}{{D}}_{1} $ 中,和平面 $ {{A}}_{1}{{DB}} $ 平行的侧面对角线有\blkda{3}条；分别是\blkda[4]{$CD_1,CB_1,D_1B_1$}.
    
    \item 在正方体 $ {{ABCD}} - {{A}}_{1}{{\;B}}_{1}{{C}}_{1}{{D}}_{1} $ 中, $ {E} $ 为 $ {{{DD}}}_{1} $ 的中点,则 $ {{{BD}}}_{1} $ 与平面 $ {{AEC}} $ 的位置关系为\blkda{平行}.
    
    \item 过直线外一点和这条直线平行的平面有\blkda{无数}个.

    \item 在空间四边形 $ {ABCD} $ 中, $ E $ 、 $ F $ 分别为 $ {AB} $ 、 $ {AD} $ 上的点,且 $ {{AE}} : {{EB}} = {{AF}} : {{FD}} = 1 : 4 $ ,又 $ H\text{、}G $ 分别为 $ {BC}\text{、}{CD} $ 的中点,则\blkda{（2）}(填序号).
    \begin{enumerate}
        \item  $ {{BD}}// $ 平面 $ {EFG} $ ,且四边形 $ {EFGH} $ 是平行四边形

        \item $ {EF}// $ 平面 $ {BCD} $ ,且四边形 $ {EFGH} $ 是梯形

        \item $ {{HG}}// $ 平面 $ {ABD} $ ,且四边形 $ {EFGH} $ 是平行四边形

        \item $ {{EH}}// $ 平面 $ {ADC} $ ,且四边形 $ {EFGH} $ 是梯形
        
    \end{enumerate}

    \begin{center}

        \includegraphics[width=0.25\linewidth]{fsajhsdfa'hsud.png}
    \end{center}

    
    \item 若直线 $ a//b $ , $ a// $ 平面 $ \alpha $ ,则 $ b $ 与平面 $ \alpha $ 的位置关系是\blkda[4]{$b//\alpha \text{或} b\subset \alpha$}.

    \item 已知直线 $ a// $ 平面 $ a $ ,点 $ A \in  a $ ,点 $ A \in $ 直线 $ b $ ,且 $ a//b $ ,则直线 $ b $ 与平面 $ a $ 的位置关系是\blkda[2]{$b\subset \alpha$}.

    
    \item 设 $ m,n $ 是平面 $ a $ 外的两条直线,给出三个论断： $ \left( 1\right) m//n,\left( 2\right) m//a,\left( 3\right) n//a $ ,以其中两个为条件, 余下的一个为结论, 写出你认为正确的一个命题：\blkda[6]{(1)(2)推(3);(1)(3)推(2)}.

    
    \item 已知 $ \bigtriangleup  {ABC} $ 中, $ {AB} = 5,{AC} = 7,A = {60}^{ \circ  } $ , $ G $ 是重心,若过 $ G $ 的平面 $ a $ 与 $ {BC} $ 平行, $ {AB} \cap  a = M,{AC} \cap  a = N $ ,则 $ {MN} = $ \blkda{$\dfrac{2\sqrt{39}}{3}$}.

    
    \item 直线与平面平行的充要条件是 \dotfill{(\kh{C})}
    \xx{直线与平面内的一条直线平行} {直线与平面内的两条直线不相交}
    {直线与平面内的任一直线都不相交} {直线与平行内的无数条直线平行}

    \item 下列命题正确的是 \dotfill{(\kh{C})}
    \xx{一直线与平面平行,则它与平面内任一直线平行}{一直线与平面平行,则平面内有且只有一个直线与已知直线平行}{ 一直线与平面平行,则平面内有无数直线与已知直线平行,它们在平面内彼此平行}{一直线与平面平行,则平面内任意直线都与已知直线异面}


    \item 若直线 $ a $ 在平面 $ \alpha $ 内,直线 $ a $ 、 $ b $ 是异面直线,则直线 $ b $ 和 $ \alpha $ 平面的位置关系是 \dotfill{(\kh{C})}
    \xx{相交} {平行} {相交或平行} {相交且垂直}

    \item 下列各命题中假命题的个数为\dotfill{(\kh{B})}
    \begin{enumerate}
        \item 经过两条平行直线中一条直线的平面必平行于另一条直线；
        \item 若一条直线平行于两相交平面,则这条直线和交线平行；
        \item 空间四边形中三条边的中点所确定平面和这个空间四边形的两条对角线都
平行.
    \end{enumerate}
    \xx{0}{1}{2}{3}

    \item 已知 $ a $ 、 $ b $ 为两异面直线,下列结论正确的是\dotfill{(\kh{D})}
    \xx{过不在 $ a\text{、}b $ 上的任何一点,可作一个平面与 $ a\text{、}b $ 都平行}
    {过不在 $ a\text{、}b $ 上的任一点,可作一直线与 $ a\text{、}b $ 都相交}
    {过不在 $ a\text{、}b $ 上任一点,可作一直线与 $ a\text{、}b $ 都平行}
    {过 $ a $ 可以并且只可以作一个平面与 $ b $ 平行}

    
    \item 如图,已知 $ {{ABCD}} $ 与 $ {{BAFE}} $ 是两个不共面的正方形,且 $ {{AM}} = {{FN}} $ .求证: $ {{MN}}// $ 平面 $ {{BEC}} $ .
    
    \begin{center}

        \includegraphics[width=0.25\linewidth]{asdfhasldf.png}
    \end{center}
    \begin{jieda}

     \begin{center}
    
            \includegraphics[width=0.25\linewidth]{bodyxb//src/fa;skjdhfa'hfa.png}
        \end{center}   
    \end{jieda}
    
    \begin{jieda}
        过M作$MG//AB$，交$BC$于点G，过N作$NK//AB$，交$BE$于点G，可以证明$GK\subset \alpha,MN\not\subset,GK//MN$,故$MN//平面BEC$.
    \end{jieda}

    
    \item 用平行于空间四边形 $ {{ABCD}} $ 的一组对边 $ {{AB}}\text{、}{{CD}} $ 的平面截此空间四边形。
    \begin{enumerate}
        \item 求证：所得截面 $ {{MNPQ}} $ 是平行四边形；
        \item 如果 $ {{AB}} = {{CD}} = a $ ,求证：四边形 $ {{MNPQ}} $ 的周长为定值.
    \end{enumerate}
  
    \begin{center}

        \includegraphics[width=0.25\linewidth]{cjshf'kashdfws.png}
    \end{center}
    
    \begin{jieda}
    \begin{enumerate}
        \item AB//平面MNPQ，$AB\subset 平面ABC, 平面ABC\cap 平面MNPQ=MN$，由线面平行的判定定理可知，$AB//MN$，同理可得$AB//PQ, CD//MQ, CD//NP$，因此$PQ//MN,NP//MQ$，因此为平行四边形.
        \item 因为$AB//MN,MQ//CD$，因此$\dfrac{BN}{NC}=\dfrac{AM}{MC}=\dfrac{AQ}{QD}=\dfrac{1-t}{t}$，从而周长$=2(ta+(1-t)a)=2a$.
    \end{enumerate}

    \end{jieda}

    
\end{enumerate}


    





\subsection*{直线与平面的位置关系（2）}

知识点：线面垂直

A 组：

\begin{enumerate}

\item 下列关于直线与平面垂直的判断中，正确的是\dotfill{(\kh{C})}

\xx{若直线 ${a}$ 与平面 ${\alpha }$ 内的一条直线垂直，则直线 ${a}$ 与平面 ${\alpha }$ 垂直}
{若直线 ${a}$ 与平面 ${\alpha }$ 内的两条平行直线垂直，则直线 ${a}$ 与平面 ${\alpha }$ 垂直}
{若直线 ${a}$ 与平面 ${\alpha }$ 内的两条相交直线垂直，则直线 ${a}$ 与平面 ${\alpha }$ 垂直}
{若直线 ${a}$ 与平面 ${\alpha }$ 内的无数条直线垂直，则直线 ${a}$ 与平面 ${\alpha }$ 垂直}

\item 给出下列命题：（1）若直线 ${l}$ 与平面 ${\alpha }$ 中的无数条直线垂直，则 ${l}\bot {\alpha }$ ；（2）若直线 ${l}\bot$ 平面 ${\alpha }$ ,且直线 ${m} \bot$ 平面 ${\alpha }$ ,则 ${l}//{m}$ ; (3) 若 ${A} \in  {l},{B} \in  {l}$ 且 ${A} \in  {\alpha },{B} \in  {\alpha }$ ,可得 ${l} \subset  {\alpha }$ . 其中真命题的个数是 $\;$ \dotfill{(\kh{C})}
\xx{0}{1}{2}{3}

\item m，n为空间两条不重合直线，α为空间平面，下列命题正确的是\dotfill{(\kh{B})}

\xx{${m} \bot  {\alpha },{n} \bot  {m}$ ,则 ${n}\parallel {\alpha }$}
{${m},{n}$ 与 ${\alpha }$ 所成角均为 ${90}^{ \circ  }$ ,则 ${m}\parallel {n}$}
{${m}\parallel {\alpha },{n}\parallel {\alpha },{m}\parallel {n}$ ,则直线 ${m},{n}$ 到 ${\alpha }$ 的距离相等}
{${m}\parallel {\alpha },{n}\parallel {\alpha }$ ,则 ${m}\parallel {n}$}

\item 下列四个正方体图形中, $l$ 是正方体的一条对角线,点 ${D}\text{ 、 }{E}\text{ 、 }{F}$ 分别为其所在棱的中点, 能得出 ${I} \bot$ 平面 ${DEF}$ 的是 \dotfill{(\kh{C})}

\begin{center}
    \centering
    \includegraphics[width=0.9\linewidth]{bodyxb//src/aouihfiaqgefiq'ihwerfq.png}
\end{center}

\item 如图，已知正方体 ${ABCD} - {A}_{1}{B}_{1}{C}_{1}{D}_{1}$ 的棱长为 1 .

\begin{center}
    \centering
    \includegraphics[width=0.25\linewidth]{bodyxb//src/akuihfai'bf'ihqwiherfvq.png}
\end{center}
\begin{enumerate}
    \item 点 $A$ 到平面 $B{B}_{1}{C}_{1}C$ 的距离为\blkda{1}；
    \item 直线 ${B}_{1}{D}_{1}$ 和平面 ${ABCD}$ 的距离为\blkda{1}；
    \item 直线 ${A}_{1}{B}_{1}$ 和平面 ${AB}{C}_{1}{D}_{1}$ 的距离为\blkda{$\dfrac{\sqrt{2}}{2}$}.
\end{enumerate}

\end{enumerate}

B 组:
\begin{enumerate}

\item 下列命题中正确的是 \dotfill{(\kh{D})}

\xx{若直线 $a \bot  a$ 内的一条直线,则 $a \bot  a$}
{若直线 $a \bot  a$ 内的无数条直线,则 $a \bot  a$}
{若 $a$ 与 $a$ 不垂直,则 $a$ 内不存在与 $a$ 垂直的直线}
{若 $a$ 与 $a$ 不垂直，则 $a$ 内不存在与 $a$ 垂直的两条相交直线}

\item 下列命题中，假命题是 \dotfill{(\kh{D})}

\xx{如果一条直线与一个平面垂直，则这条直线与这个平面内的任意一条直线垂直}
{如果两条直线同垂直于一个平面，则这两条直线平行}
{平面的垂线与这个平面一定相交}
{如果一条直线上有两点到一个平面的距离相等，则这条直线与这个平面平行}

\item 下列四个命题，其中正确命题的个数是 \dotfill{(\kh{C})}
\begin{enumerate}
    \item 点 $A \notin$ 平面 $\alpha$ ，直线 $a \subset  a$ ，则 $A$ 到 $a$ 的距离等于 $A$ 到 $\alpha$ 的距离；
    
    \item 直线 $a\parallel$ 平面 $\alpha$ ， $b \subset  a$ ， $a\parallel b$ ，则 $a$ 、 $b$ 间的距离等于 $a$ 与 $\alpha$ 间的距离；
    
    \item $A \in$ 直线 $a,a.m$ 直线 $b$ ，则 $a\text{ 、 }b$ 间的距离等于 $A$ 到 $b$ 的距离；
        
    \item 直线 $a\parallel$ 平面 $\alpha ,A \in  a$ ，则 $a\text{ 、 }\alpha$ 间的距离等于 $A$ 到 $\alpha$ 的距离。
\end{enumerate}

\xx{4 个} {3 个} {2 个} {1 个}

\item 从平面 $a$ 外一点 P 作与 $a$ 相交的直线，使 P 与交点的距离为 1，则满足条件的直线不可能是\dotfill{(\kh{C})}

\xx{0 条} {1 条} {2 条} {无数条}

\item 若正 $\bigtriangleup {ABC}$ 的边长为 $\dfrac{4\sqrt{3}}{3}$ ，则到三个顶点的距离都为 1 的平面有 \dotfill{(\kh{C})}

\xx{1 个} {3 个} {5 个} {7 个}

\item “直线 $l$ 垂直于平面 $\alpha$ 内的无数条直线”是“ $l\bot a$ ”的\blkda{必要非充分}条件.

\item 已知 $m,n$ 是空间两条相交直线， ${l}_{1},{l}_{2}$ 是与 $m,n$ 都垂直的两条直线，直线 $l$ 与 ${l}_{1},{l}_{2}$ 都相交，则直线 $l$ 与 ${l}_{1},{l}_{2}$ 所成的角的大小关系为\blkda{相等}.

\item 已知三条相交于一点的线段 ${PA}\text{ 、 }{PB}\text{ 、 }{PC}$ 两两垂直，且 $A\text{ 、 }B\text{ 、 }C$ 在同一平面内， $P$ 在平面 ${ABC}$ 外， ${PH}\bot$ 平面 ${ABC}$ 于 $H$ ，则垂足 $H$ 是 $\bigtriangleup  {ABC}$ 的\blkda{垂}心。


\item 已知BC是Rt $\bigtriangleup  {ABC}$ 的斜边， ${AP}\bot$ 平面 ${ABC}$ ，连结 ${PB}\text{ 、 }{PC}$ ，作 ${PD}\bot {BC}$ 于 $D$ ，连结 ${AD}$ ， 则图中共有直角三角形\blkda{8}个， $\bigtriangleup  {PBC}$ 是\blkda{锐角}三角形.

\item 在棱长为 $\alpha$ 的正方体 ${ABCD} - {A}_{1}{B}_{1}{C}_{1}{D}_{1}$ 中， $A$ 到直线 ${B}_{1}C$ 的距离为\blkda{$\dfrac{\sqrt{6}a}{2}$}， $A{A}_{1}$ 到平面 $B{B}_{1}{D}_{1}D$ 的距离为\blkda{$\dfrac{\sqrt{2}a}{2}$}， $A$ 到平面 ${A}_{1}{BD}$ 的距离为\blkda{$\dfrac{\sqrt{3}a}{3}$}.


\item 已知线段 ${AB}$ 的两端点 $A$ 、 $B$ 到平面 $\alpha$ 的距离分别为 $a$ 和 $b$ ， （1）若 $A$ 、 $B$ 在 $\alpha$ 的同侧，则 ${AB}$ 中点到 $\alpha$ 的距离为\blkda{$\dfrac{a+b}{2}$}.

（2）若 $A$ 、 $B$ 在 $\alpha$ 的异侧，则 ${AB}$ 中点到 $\alpha$ 的距离为\blkda{$\dfrac{|a-b|}{2}$}.

\item 若三角形 ${ABC}$ 三个顶点到平面 $\alpha$ 的距离分别为 $a$ 、 $b$ 、 $c$ ， $\bigtriangleup {ABC}$ 的重心为 $G$ ，则点 $G$ 到平面 $\alpha$ 的距离为\blkda[10]{$\dfrac{a+b+c}{3}$,$\dfrac{|-a+b+c|}{3}$,$\dfrac{|a-b+c|}{3}$,$\dfrac{|a+b-c|}{3}$}.

\item 如图 ABCD 是矩形， ${PA} \bot$ 平面 ${ABCD}$ ， $\bigtriangleup  {PAD}$ 是等腰三角形，M，N 分别是 AB，PC 的中点,求证: ${MN} \bot$ 平面 ${PCD}$ .

\begin{center}
    \centering
    \includegraphics[width=0.25\linewidth]{bodyxb//src/aufh'fibqirvbf'qiyefq.png}
\end{center}

\begin{jieda}
    取 \( {PD} \) 的中点 \( E \) ,连接
\( {AE},{NE} \) ,
\( \because N \) 为 \( {PC} \) 的中点,
\( \therefore {NE} \) 为 \( \bigtriangleup  {PCD} \) 的中位线,
\( \therefore {NE}\overset{//}{ = }\dfrac{1}{2}{CD} \) .
在矩形 \( {ABCD} \) 中, \( {AB} ⫫ {CD} \) ,
又 \( M \) 为 \( {AB} \) 的中点, \( \therefore {AM}\overset{//}{ = }\dfrac{1}{2}{CD} \) .
\( \therefore {AM}\overset{//}{ = }{NE} \) ,
\( \therefore \) 四边形 \( {AMNE} \) 为平行四边形, \( \therefore {AE}//{MN} \) .
\( \because {PA} \bot \) 平面 \( {ABCD},{AD} \subset \) 平面 \( {ABCD} \) ,
\( \therefore {PA}\bot {AD} \) .
又 \( \bigtriangleup {PAD} \) 为等腰三角形, \( E \) 为 \( {PD} \) 的中点,
\( \therefore {PA} = {AD},{AE} \bot  {PD},\therefore {MN} \bot  {PD} \) .
\( \because {PA}\bot \) 平面 \( {ABCD},{CD} \subset \) 平面 \( {ABCD} \) ,
\( \therefore {PA}\bot {CD} \) .
在矩形 \( {ABCD} \) 中, \( {CD} \bot  {AD} \) ,
又 \( {PA} \cap  {AD} = A,{PA},{AD} \subset \) 平面 \( {PAD} \) ,
\( \therefore {CD} \bot \) 平面 \( {PAD} \) ,
\( \because {AE} \subset \) 平面 \( {PAD},\therefore {AE} \bot  {CD},\therefore {MN} \bot  {CD} \) ,
又 \( {PD} \cap  {CD} = D,{PD},{CD} \subset \) 平面 \( {PCD} \) ,
\( \therefore {MN}\bot \) 平面 \( {PCD} \) .
\end{jieda}

\item 如图，已知 ${ABCD}$ 是矩形， ${AB} = a,{AD} = b$ ， ${PA} \bot$ 平面 ${ABCD}$ ， ${PA} = {2c}$ ， $Q$ 是 ${PA}$ 的中点. 求
\begin{enumerate}
    \item $Q$ 到 ${BD}$ 的距离；
    \item $P$ 到平面 ${BQD}$ 的距离.
\end{enumerate}

\begin{center}
    \centering
    \includegraphics[width=0.25\linewidth]{bodyxb//src/afuhg'igq'iwyf'yiq.png}
\end{center}

\begin{jieda}
    \begin{enumerate}
        \item 在矩形 \( {ABCD} \) 中,作 \( {AE}\bot {BD},E \) 为垂足,连
结 \( {QE} \) .
\( \because {QA} \bot \) 平面 \( {ABCD} \) ,由三垂线定理得 \( {QE} \bot  {BE} \) ,
\( \therefore {QE} \) 的长为 \( Q \) 到 \( {BD} \) 的距离.
在矩形 \( {ABCD} \) 中, \( {AB} = a,{AD} = b \) ,
\( \therefore {AE} = \dfrac{ab}{\sqrt{{a}^{2} + {b}^{2}}} \) , 图 3-2-101
在 Rt \( \bigtriangleup {QAE} \) 中, \( {QA} = \dfrac{1}{2}{PA} = c \) ,
\( \therefore {QE} = \sqrt{Q{A}^{2} + A{E}^{2}} = \sqrt{{c}^{2} + \dfrac{{a}^{2}{b}^{2}}{{a}^{2} + {b}^{2}}}\therefore Q \) 到 \( {BD} \) 的距离为 \( \sqrt{{c}^{2} + \dfrac{{a}^{2}{b}^{2}}{{a}^{2} + {b}^{2}}} \) .
    \item 解法一 因为平面 \( {BQD} \) 经过线段 \( {PA} \) 的中点,
\( \therefore P \) 到平面 \( {BQD} \) 的距离等于 \( A \) 到平面 \( {BQD} \) 的距离.
在 \( \bigtriangleup {AQE} \) 中, \( {AH}\bot {QE},H \) 为垂足.
\( \because {BD} \bot  {AE},{BD} \bot  {QE},\therefore {BD} \bot \) 平面 \( {AQE}.\therefore {BD} \bot  {AH} \) .
\( \therefore {AH} \bot \) 平面 \( {BQE} \) ,即 \( {AH} \) 为 \( A \) 到平面 \( {BQD} \) 的距离.
在 Rt \( \bigtriangleup {AQE} \) 中, \( \because {AQ} = c,{AE} = \dfrac{ab}{\sqrt{{a}^{2} + {b}^{2}}} \) ,
\( \therefore {AH} = \dfrac{abc}{\sqrt{\left( {{a}^{2} + {b}^{2}}\right) {c}^{2} + {a}^{2}{b}^{2}}} \) .
\( \therefore P \) 到平面 \( {BD} \) 的距离为 \( \dfrac{abc}{\sqrt{\left( {{a}^{2} + {b}^{2}}\right) {c}^{2} + {a}^{2}{b}^{2}}} \) .\\
解法二 本题也可用体积法求解
设点 \( A \) 到平面 \( {QBD} \) 的距离为 \( h \) ,
由 \( {V}_{A - {BQD}} = {V}_{Q - {ABD}} \) 得
\( \dfrac{1}{3}{S}_{\bigtriangleup {BQD}} \cdot  h = \dfrac{1}{3}{S}_{\bigtriangleup {ABD}} \cdot  {AQ} \)
\( h = \dfrac{{S}_{\bigtriangleup {ABD}} \cdot  {AQ}}{{S}_{\bigtriangleup {BQD}}} = \cdots  = \dfrac{abc}{\sqrt{\left( {{a}^{2} + {b}^{2}}\right) {c}^{2} + {a}^{2}{b}^{2}}} \)
    \end{enumerate}
\end{jieda}



\end{enumerate}

C 组:
\begin{enumerate}

\item 如图，在边长为 1 的正方体 ${{ABCD}} - {{A}}_{1}{{B}}_{1}{{C}}_{1}{{D}}_{1}$ 中， ${E},{F},{M}$ ， ${N}$ 分别为棱 ${{AB}}$ ， ${AD},{D}_{1}{C}_{1},{C}_{1}{B}_{1}$ 的点,满足 $\dfrac{{C}_{1}M}{{C}_{1}{D}_{1}} = \dfrac{{C}_{1}N}{{C}_{1}{B}_{1}} = \dfrac{AE}{AB} = \dfrac{AF}{AD} = \lambda \left( {\lambda  \in  \left( {0,1}\right) }\right)$ ,过 $E,F$ , ${M},{N}$ 四点作该正方体的截面，则下列说法错误的是 \dotfill{(\kh{BD})}

\begin{center}
    \centering
    \includegraphics[width=0.25\linewidth]{bodyxb//src/oafhqouh3rofhow3uf.png}
\end{center}

\xx{$\lambda  = \dfrac{7}{2}$ 时，该截面是正六边形}
{$\lambda  = \dfrac{7}{8}$ 时，四边形 ${E}{F}{M}{N}$ 为正方形}
{${MN}//$ 平面 $A{D}_{1}{B}_{1}$}
{当四边形${E}{F}{M}{N}$ 为正方形时,它的面积为 $\dfrac{8}{9}$}

\item 如图，在平行四边形 ${{ABCD}}$ 中， ${{BC}} = 2{{AB}} = {4}$ ， $\angle {{ABC}} = {{{60}}}^{ \circ  }$ ， ${{PA}} \bot$ 平面 ${{ABCD}},{{PA}} = 2,{E},{F}$ 分别为 ${{BC}},{{PE}}$ 的中点,求证: ${{AF}} \bot$ 平面 ${{PED}}$ .

\begin{center}
    \centering
    \includegraphics[width=0.25\linewidth]{bodyxb//src/gsh'eouthgouergsbje.png}
\end{center}

\begin{jieda}
    证明：如图,连接 \( {AE} \) ,在平行四边形 \( {ABCD} \) 中,
\( {BC} = {2AB} = 4 \) ,
\( \angle {ABC} = {60}^{ \circ  },E \) 是 \( {BC} \) 的中点,
所以 \( {AE} = 2,{ED} = 2\sqrt{3} \) ,
从而有 \( A{E}^{2} + E{D}^{2} = A{D}^{2} \) , 所以 \( {AE}\bot {ED} \) ,
又 \( {PA} \bot \) 平面 \( {ABCD},{ED} \subset \) 平面 \( B \) \( C \)
\( {ABCD} \) ,
所以 \( {PA}\bot {ED} \) ,
因为 \( {PA} \cap  {AE} = A,{PA},{AE} \subset \) 平面 \( {PAE} \) ,
所以 \( {ED}\bot \) 平面 \( {PAE} \) ,
又 \( {AF} \subset \) 平面 \( {PAE} \) ,所以 \( {ED} \bot  {AF} \) .
又 \( {PA} = {AE} = 2,F \) 为 \( {PE} \) 的中点,
所以 \( {AF}\bot {PE} \) ,
又 \( {PE} \cap  {ED} = E,{PE},{ED} \subset \) 平面 \( {PED} \) ,
所以 \( {AF} \bot \) 平面 \( {PED} \) .


\end{jieda}

\item 如图， ${{PA}} \bot$ 矩形 ${ABCD}$ ，有下列结论：① ${{PB}} \bot  {{BC}}$ ， ${{QPD}} \bot  {{CD}}$ ，③ ${{PD}} \bot  {BD}$ ， ④ ${{PA}} \bot  {{BD}}$ . 其中正确的是\blkda{(1)(2)(4)}(填序号).
\begin{center}
    \centering
    \includegraphics[width=0.25\linewidth]{bodyxb//src/lajerfoawbfojaer.png}
\end{center}


\end{enumerate}



\clearpage

\subsection*{直线与平面的位置关系（3）}

知识点：线面所成角

A组：
\begin{enumerate}
\item 直线 ${l}$ 是平面 ${\alpha }$ 的一条斜线， ${l}$ 与平面 ${\alpha }$ 内的直线所成角的取值范围是\blkda{$[<l,\alpha>,\dfrac{\pi}{2}]$}.

\item 已知空间中两条直线 ${a},{b}$ 无公共点，则 “直线 ${a},{b}$ 与平面 ${\alpha }$ 所成的角相等” 是 “ ${a}//{b}$ ” 的\dotfill{(\kh{B})}
\xx{充分不必要条件}{必要不充分条件}{充要条件}{既不充分也不必要条件}

\item 若直线 ${l}$ 与平面 ${\alpha }$ 所成的角为 ${\theta }$ ，直线 ${l}$ 与平面 ${\alpha }$ 上的直线 ${m}$ 所成的角为 ${\varphi }$ ，则总有 \dotfill{(\kh{B})}
\xx{${\theta } < {\varphi }$ }{${\theta } \leq  {\varphi }$}{C. ${\theta } > {\varphi }$ }{D. ${\theta } \geq  {\varphi }$}

\item 直线 $l$ 与平面 ${\alpha }$ 相交， $A,B$ 在 $l$ 上， ${{AB}} = {4},{{AB}}$ 在 ${\alpha }$ 内的射影长为 $2\sqrt{2}$ ， ${l}$ 与 ${\alpha }$ 夹角为\blkda{$\dfrac{\pi}{4}$}.

\item 在正方体 ${{ABCD}} - {{A}}_{1}{{B}}_{1}{{C}}_{1}{{D}}_{1}$ 中， $E$ 是 ${A}{{A}}_{1}$ 的中点，则直线 ${{B}}_{1}{E}$ 与平面 BCE所成角的正弦值为\blkda{$\dfrac{4}{5}$}.
    
\end{enumerate}


B 组:
\begin{enumerate}


\item 下列命题：\ding{172}一条直线在平面内的射影是一条直线；\ding{173} ${Rt} \bigtriangleup  {ABC}$ 的斜边在平面内，直角顶点 $C$ 在平面内的射影为 ${C}^{\prime }$ ，则 $\bigtriangleup  {AB}{C}^{\prime }$ 是钝角三角形；\ding{174}在同一平面内的射影长相等，则斜线长相等；\ding{175}两斜线与平面所成的角相等，则这两斜线互相平行.其中真命题的个数是\blkda{0}.

\item 下列命题中正确的是\dotfill{(\kh{A})}

\xx{若平面 $M$ 外的两条直线在平面 $M$ 内的射影为一条直线及此直线外的一个点, 则它们互为异面直线}
{若平面 $M$ 外的两条直线在平面 $M$ 内的射影为两条相交直线，则这两条直线相交}
{若平面 $M$ 外的两条直线在平面 $M$ 内的射影为两条平行直线，则这两条直线平行}
{若平面 $M$ 外的两条直线在平面 $M$ 内的射影为两条互相垂直的直线，则这两条直线垂直}

\item 已知直线 $a$ 与平面 $\alpha$ 所成的角为 ${30}^{ \circ  },b \subset  a$ 且与 $a$ 异面，若直线 $a$ 与 $b$ 所成的角为 $\varphi$ ,则\dotfill{(\kh{C})}

\xx{${0}^{ \circ  } < \varphi  \leq  {30}^{ \circ  }$ }{ ${0}^{ \circ  } < \varphi  \leq  {90}^{ \circ  }$}{  ${30}^{ \circ  } \leq  \varphi  \leq  {90}^{ \circ  }$}{ ${30}^{ \circ  } \leq  \varphi  \leq  {180}^{ \circ  }$}


\item 已知 ${AB}\bot$ 平面 $a$ 于 $B$ ， ${BC}$ 为 ${AC}$ 在 $a$ 内的射影， ${CD} \subset  a$ ，若 $\angle {ACD} = {60}^{ \circ  }$ ， $\angle {BCD} = {45}^{ \circ  }$ ，则 ${AC}$ 与平面 $a$ 所成的角为\blkda{$\dfrac{\pi}{4}$}.

\item 在正方体 ${ABCD} - {A}_{1}{B}_{1}{C}_{1}{D}_{1}$ 中，直线 $A{B}_{1}$ 与平面 ${AB}{C}_{1}{D}_{1}$ 所成的角的大小为\blkda{$\dfrac{\pi}{6}$}.

\item 已知 $P$ 是 $\bigtriangleup  {ABC}$ 所在平面外一点， $O$ 是 $P$ 点在平面 $\alpha$ 上的射影。若 $P$ 到 $\bigtriangleup  {ABC}$ 三边的距离相等，则 $O$ 是 $\bigtriangleup  {ABC}$ 的\blkda{内心}，若 $P$ 到 $\bigtriangleup  {ABC}$ 三个顶点的距离相等，则 $O$ 是 $\bigtriangleup  {ABC}$ 的\blkda{外心}，若 ${PA}\text{ 、 }{PB}\text{ 、 }{PC}$ 两两互相垂直，则 $O$ 是 $\bigtriangleup  {ABC}$ 的\blkda{垂心}，

\item 已知 $\bigtriangleup  {ABC}$ 中， $A \in  a$ ， ${BC}//a$ ， ${BC} = 6$ ， $\angle {BAC} = {90}^{ \circ  }$ ， ${AB}$ 、 ${AC}$ 与平面 $a$ 分别成 ${30}^{ \circ  }$ 、 ${45}^{ \circ  }$ 的角。则 ${BC}$ 到平面 $\alpha$ 的距离为\blkda{$\sqrt{6}$}.

\item 已知直角 $\bigtriangleup  {ABC}$ 的斜边 ${AB}$ 在平面 $M$ 内，而直角边 ${AC}$ ， ${BC}$ 分别与 $M$ 成 ${30}^{ \circ  }$ ， ${45}^{ \circ  }$ 角，则斜边上的高 ${CD}$ 与平面 $M$ 所成的角为\blkda{$\dfrac{\pi}{3}$}.

\item 已知 ${PA}\text{ 、 }{PB}\text{ 、 }{PC}$ 是从点 $P$ 发出的不共面的三条射线，每两条射线的夹角都是 ${60}^{ \circ  }$ ， 则直线 PC 与平面 ${PAB}$ 所成的角的余弦值为\blkda{$\dfrac{\sqrt{3}}{3}$}.

\item 到空间不共面的四点距离相等的平面有\blkda{7}个.

\item 已知长为 3 的线段 ${AB}$ 的两个端点到平面 ${\alpha }$ 的距离分别为 1 和 2，则直线 ${AB}$ 与平面 ${\alpha }$ 的所成角大小为\blkda{$arcsin\dfrac{1}{3}or\dfrac{\pi}{2}$}.

\item 已知 ${\theta }$ 是锐角，则 “直线 ${l}$ 与平面 ${\alpha }$ 所成角的大小为 ${\theta }$ ” 是 “直线 ${l}$ 与平面 ${\alpha }$ 内无数条直线所成角的大小为 ${\theta }$ ” 的\dotfill{(\kh{B})}条件.

\xx{必要不充分 }{充分不必要}
{充分必要} {既不充分也不必要}


\item 如图，已知正方体 ${ABCD} - {EFGH}$ 的棱长为 $a$ ，点 $P$ 在 ${AC}$ 上， $Q$ 在 ${BG}$ 上， ${AP} = {BQ} = a$ . 求直线 ${PQ}$ 与平面 ${ABCD}$ 所成角的大小.

\begin{center}
    \centering
    \includegraphics[width=0.25\linewidth]{bodyxb//src/hfa'suodhf'os.png}
\end{center}

\begin{jieda}
    过点$Q$作$BC$的垂线，交BC于点$E$,可证明$E$即为$Q$在平面$ABCD$的投影，因此 $\angle{QPE}$即为所求角，$|QE|=\dfrac{\sqrt{2}}{2}a,|EP|=(1-\dfrac{\sqrt{2}}{2})a$,$tan\angle{QPE}=\sqrt{2}+1$，综上所求角为$arctan(\sqrt{2}+1)$

    
    \begin{center}

        \includegraphics[width=0.25\linewidth]{bodyxb//src/ashfaohfaoijf.png}
    \end{center}
    
\end{jieda}

\item 如图，在 ${Rt\Delta BMC}$ 中，斜边 ${BM} = 5$ ，它在平面 ${ABC}$ 上的射影 $AB$ 长为 4， $\angle {MBC} = {60}^{ \circ  }$ ,求 MC 与平面 MAB 所成角的大小.

\begin{center}
    \centering
    \includegraphics[width=0.25\linewidth]{bodyxb//src/ahfa'hofuafa.png}
\end{center}

\begin{jieda}
    过C作AB的垂线，垂足为N，因为$CN\perp AB$,$CN\perp AM$，因此$CN\perp 面ABM$,所以$\angle{NMC}$为所成角. 因为$BC\perp CM$，$BC\perp AM$，因此$BC\perp 面ACM$,所以$BC\perp AC$，由等面积法$S_{\triangle ABC}=\dfrac{1}{2}|BC||AC|=\dfrac{1}{2}|AB||CN|$，求得$|CN|=\dfrac{5\sqrt{39}}{16}$. 又因为$|CM|=\dfrac{5\sqrt{3}}{2}$,所以$sin\angle{NMC}=\dfrac{CN}{CM}=\dfrac{\dfrac{5\sqrt{39}}{16}}{\dfrac{5\sqrt{3}}{2}}=\dfrac{\sqrt{13}}{8}$。综上，所求角为$arcsin \dfrac{\sqrt{13}}{8}$.
    \begin{center}
        \centering
        \includegraphics[width=0.25\linewidth]{bodyxb//src/fasuhfaohfa.png}
    \end{center}
\end{jieda}



\item 如图：在平面 $\beta$ 内有 $\bigtriangleup  {ABC}$ ，在平面 $\beta$ 外有点 $S$ ，斜线 ${SA}\bot {AC}$ ， ${SB}\bot {BC}$ ， 且斜线 ${SA}\text{ 、 }{SB}$ 分别与平面 $\beta$ 所成的角相等.
\begin{enumerate}
    \item 求证： ${AC} = {BC}$ ；
    \item 又设点 $S$ 与平面 $\beta$ 的距离是 ${4cm}$ ， ${AC}\bot {BC}$ ，且 ${AB} = {6cm}$ ， 求 ${SC}$ 与平面 $\beta$ 所成的角.
\end{enumerate}

\begin{center}
    \centering
    \includegraphics[width=0.5\linewidth]{bodyxb//src/faushf'ahfoaihjfa.png}
\end{center}

\begin{jieda}
    \begin{enumerate}
        \item 设S在$\beta$的投影为D，因为$\angle{SAD}=\angle{SBD}$，所以$SA=SB$,又因为$\angle{SAC}=\angle{SBC}$,所以AC=BC.
        \item $AC\perp面SAC$,所以$AC\perp AD$,同理$CB\perp BD$，又因为AC=BC,所以四边形ACBD为正方形. CD=AB=6,SD=4,所以$tan\angle{SCD}=\dfrac{4}{6}=\dfrac{2}{3}$.综上，所成角为$arctan\dfrac{2}{3}$
    \end{enumerate}
\end{jieda}

\item 在空间四边形 ${ABCD}$ 中， ${AB} = {CD} = 8$ ， $E$ 、 $F$ 分别为 ${AC}$ 、 ${BD}$ 的中点.

\begin{enumerate}
    \item 若 ${EF} = {4\sqrt{3}}$ ，求 ${AB}$ 与 ${CD}$ 所成的角；
    \item 若 ${AB}$ 与 ${CD}$ 所成的角为 ${60}^{ \circ  }$ ， 求 EF 的长.
\end{enumerate}

\begin{jieda}
    \begin{enumerate}
        \item 取BC的中点G，根据余弦定理，设所成角为$\theta$，有$cos\theta=|cos\angle{EGF}|=|\dfrac{4^2+4^2-(4\sqrt{3})^2}{2*4*4}|=\dfrac{1}{2}$，所以所成角为$\dfrac{\pi}{3}$.
        \item 当$\angle{EGF=\dfrac{\pi}{3}}$时，$EF=4$;当$\angle{EGF=\dfrac{2\pi}{3}}$时，$EF=4\sqrt{3}$.
    \end{enumerate}
\end{jieda}

\item 在空间四边形 ${ABCD}$ 中,已知 ${AB} = {BC} = {CD} = {DA} = {AC} = {BD} = a,E,F$ 分别是 ${AB},{CD}$ 的中点. 
\begin{enumerate}
    \item 求证: ${EF}$ 和 ${AD}$ 为异面直线;
    \item 求证: ${EF}$ 是异面直线 ${AB}$ 、 ${CD}$ 的公垂线；
    \item 求异面直线 ${AB}$ 、 ${CD}$ 间的距离.
\end{enumerate}

\begin{jieda}
    \begin{enumerate}
        \item 反证法，假设E、F、A、D共面，则$B\in AE\subset \alpha$,$C\in DF\subset \alpha, A\in \alpha, D\in \alpha$,从而$AD,BC\subset \alpha$,矛盾.
        \item 在等腰三角形AFB中，E为AB的中点，所以$EF\perp AB$,同理$EF\perp CD$,且EF与AB、CD均相交，因此是公垂线.

        \item $|EF|=\dfrac{a}{2}$,$|FG|=\dfrac{a}{2}$,$\angle{EGF}=\dfrac{\pi}{2}$ ,因此$d=|EF|=\dfrac{\sqrt{2}}{2}a$
    \end{enumerate}
\end{jieda}



\item 
\begin{enumerate}
    \item 如图，已知空间四边形 ${ABCD}$ 中， ${AB} = {CD} = 3$ ， $E\text{ 、 }F$ 分别是 ${BC}\text{ 、 }{AD}$ 上的点，且 ${BE} : {EC} = {AF} : {FD} = 1 : 2$ , ${EF} = 2\sqrt{2}$ ，求 ${AB}$ 和 ${CD}$ 所成角的大小.
    
\begin{center}
    \centering
    \includegraphics[width=0.25\linewidth]{bodyxb//src/fasha'sjdhf'asfh.png}
\end{center}
    
    \item 如图，已知空间四边形 ${ABCD}$ 中各边长和对角线长 $a,E}\text{ 、 }\;F}$ 分别为边 ${AD}}\text{ 、 }{BC}}$ 的中点，求 ${AF}}$ 与 ${CE}}$ 所成的角的大小.
    
\begin{center}
    \centering
    \includegraphics[width=0.25\linewidth]{afha'ouhfaof.png}
\end{center}

    \item 在空间四边形 ${ABCD}$ 中， ${AD} = 1$ ， ${BC} = \sqrt{3}$ ， ${AD}\bot {BC}$ ，对角线 ${BD} = \dfrac{\sqrt{13}}{2}$ ， ${AC} = \dfrac{\sqrt{3}}{2}$ ，求 ${AC}$ 与 ${BD}$ 所成角的大小
    
\begin{center}
    \centering
    \includegraphics[width=0.25\linewidth]{bodyxb//src/aisufghaiuhfa.png}
\end{center}
\end{enumerate}


\begin{jieda}
    \begin{enumerate}
        \item 取BD的靠近B的三等分点，$\angle{EMF}$即为所成角或其补角，$cos\angle{EMF}=\dfrac{-3}{4}$，综上所成角为$arccos\dfrac{3}{4}$

        \item 连接FD，取FD中点为O，$\angle{OEC}$即为所求角或其补角，\newline $cos\angle{OEC}=\dfrac{{\dfrac{\sqrt{3}a}{4}}^2+{\dfrac{\sqrt{3}a}{2}}^2-{\dfrac{\sqrt{7}a}{4}}^2}{2*\dfrac{\sqrt{3}a}{4}*\dfrac{\sqrt{3}a}{2}}=\dfrac{2}{3}$，综上所求即为$arccos\dfrac{2}{3}$.

        \begin{center}
    
            \includegraphics[width=0.25\linewidth]{bodyxb//src/aofuahsoduhfaouhfa.png}
        \end{center}

        \item 取AB、AC、AD、CD的中点分别为E、F、G、H，则$\angle{EGH}$为所成角，$|EG|=\dfrac{\sqrt{13}}{4}$,$|GH|=\dfrac{\sqrt{3}}{4}$,$|EH|=\sqrt{EF^2+FH^2}=1$,所以$cos\angle{EGH}=0$，故$\dfrac{\pi}{2}$

        \begin{center}
    
            \includegraphics[width=0.25\linewidth]{asuhfouhfaoh.png}
        \end{center}
    \end{enumerate}
\end{jieda}
    
\end{enumerate}


\clearpage





